\textbf{Phần 1. Trắc nghiệm lựa chọn}
\setcounterpageref{ex}{0}

\begin{ex} %[5H9V52]
Một bể nước hình hộp chữ nhật có kích thước bên trong là: dài $2\text{ m}$, rộng $1{,}5\text{ m}$ và cao $1{,}2\text{ m}$. Bể đang chứa nước đến $60\%$ chiều cao của bể. Thể tích nước hiện có trong bể là bao nhiêu lít?
\choice
{$1\,800\text{ lít}$}
{\True $2\,160\text{ lít}$}
{$2\,400\text{ lít}$}
{$3\,600\text{ lít}$}
\loigiai{
Thể tích của bể nước là:
$$2 \times 1{,}5 \times 1{,}2 = 3{,}6\text{ (m}^3\text{)} = 3\,600\text{ (dm}^3\text{)} = 3\,600\text{ lít}.$$
Lượng nước hiện có trong bể là:
$$3\,600 \times 60\% = 2\,160\text{ (lít)}.$$
\textbf{Đáp án của câu hỏi là B}
}
\end{ex}

\begin{ex} %[5H9V52]
Một bể cá hình hộp chữ nhật có diện tích đáy là $25\text{ dm}^2$ đang chứa nước. Khi người ta thả một hòn non bộ vào trong bể thì hòn non bộ chìm hoàn toàn trong nước và mực nước trong bể dâng cao thêm $8\text{ cm}$. Thể tích của hòn non bộ đó là
\choice
{$200\text{ dm}^3$}
{\True $20\text{ dm}^3$}
{$2\text{ dm}^3$}
{$2\,000\text{ dm}^3$}
\loigiai{
Đổi: $8\text{ cm} = 0{,}8\text{ dm}$.\\
Thể tích hòn non bộ chính bằng thể tích phần nước dâng lên:\\
$$V = 25 \times 0{,}8 = 20\text{ (dm}^3\text{)}.$$
\textbf{Đáp án của câu hỏi là B}
}
\end{ex}

\begin{ex} %[5H9V52]
Một hình hộp chữ nhật có thể tích $360\text{ cm}^3$, chiều dài $10\text{ cm}$ và chiều rộng $6\text{ cm}$. Nếu tăng chiều cao của hình hộp chữ nhật đó thêm $2\text{ cm}$ thì thể tích mới của nó là
\choice
{$420\text{ cm}^3$}
{\True $480\text{ cm}^3$}
{$450\text{ cm}^3$}
{$500\text{ cm}^3$}
\loigiai{
Diện tích đáy của hình hộp chữ nhật là:
$$10 \times 6 = 60\text{ (cm}^2\text{)}.$$
Chiều cao ban đầu của hình hộp chữ nhật là:
$$360 : 60 = 6\text{ (cm)}.$$
Chiều cao mới của hình hộp là:
$$6 + 2 = 8\text{ (cm)}.$$
Thể tích mới của hình hộp chữ nhật là:
$$60 \times 8 = 480\text{ (cm}^3\text{)}.$$
\textbf{Đáp án của câu hỏi là B}
}
\end{ex}

\begin{ex} %[5H9V52]
Một khối kim loại hình hộp chữ nhật có chiều dài $15\text{ cm}$, chiều rộng $8\text{ cm}$ và chiều cao $5\text{ cm}$. Biết mỗi xăng-ti-mét khối kim loại đó nặng $7{,}8\text{ g}$. Cả khối kim loại đó cân nặng bao nhiêu ki-lô-gam?
\choice
{$3{,}9\text{ kg}$}
{\True $4{,}68\text{ kg}$}
{$46{,}8\text{ kg}$}
{$0{,}468\text{ kg}$}
\loigiai{
Thể tích khối kim loại là:
$$15 \times 8 \times 5 = 600\text{ (cm}^3\text{)}.$$
Khối lượng của khối kim loại là:
$$600 \times 7{,}8 = 4\,680\text{ (g)} = 4{,}68\text{ (kg)}.$$
\textbf{Đáp án của câu hỏi là B}
}
\end{ex}

\begin{ex} %[5H9V52]
Một bể nước hình hộp chữ nhật có chiều dài $1{,}8\text{ m}$, chiều rộng $1{,}2\text{ m}$ và chiều cao $1\text{ m}$. Người ta mở một vòi nước chảy vào bể, mỗi phút vòi chảy được $36\text{ lít}$. Biết ban đầu bể hoàn toàn cạn nước. Sau bao lâu thì nước đầy bể?
\choice
{$50\text{ phút}$}
{\True $60\text{ phút}$}
{$70\text{ phút}$}
{$80\text{ phút}$}
\loigiai{
Thể tích bể nước là:
$$1{,}8 \times 1{,}2 \times 1 = 2{,}16\text{ (m}^3\text{)} = 2\,160\text{ dm}^3 = 2\,160\text{ lít}.$$
Thời gian để nước chảy đầy bể là:
$$2\,160 : 36 = 60\text{ (phút)} = 1\text{ (giờ)}.$$
\textbf{Đáp án của câu hỏi là B}
}
\end{ex}

\begin{ex} %[5H9V52]
Người ta xếp $300$ hộp chè dạng hình hộp chữ nhật có kích thước $10\text{ cm} \times 8\text{ cm} \times 5\text{ cm}$ vào một thùng các-tông lớn. Hỏi thể tích của thùng các-tông đó ít nhất phải là bao nhiêu đề-xi-mét khối (coi thể tích các vách ngăn không đáng kể)?
\choice
{$100\text{ dm}^3$}
{\True $120\text{ dm}^3$}
{$150\text{ dm}^3$}
{$200\text{ dm}^3$}
\loigiai{
Thể tích của một hộp chè là:
$$10 \times 8 \times 5 = 400\text{ (cm}^3\text{)}.$$
Thể tích của $300$ hộp chè là:
$$400 \times 300 = 120\,000\text{ (cm}^3\text{)}.$$
Đổi $120\,000\text{ cm}^3 = 120\text{ dm}^3$.\\
\textbf{Đáp án của câu hỏi là B}
}
\end{ex}

\begin{ex} %[5H9V52]
Một đoạn mương dẫn nước hình hộp chữ nhật dài $50\text{ m}$, chiều rộng đáy mương là $1{,}2\text{ m}$ và mực nước trong mương sâu $0{,}8\text{ m}$. Lượng nước hiện có trong đoạn mương đó là
\choice
{$42\text{ m}^3$}
{\True $48\text{ m}^3$}
{$54\text{ m}^3$}
{$60\text{ m}^3$}
\loigiai{
Thể tích nước trong đoạn mương là:
$$50 \times 1{,}2 \times 0{,}8 = 48\text{ (m}^3\text{)}.$$
\textbf{Đáp án của câu hỏi là B}
}
\end{ex}

\begin{ex} %[5H9V52]
Nếu chiều dài hình hộp chữ nhật tăng gấp đôi, chiều rộng tăng gấp ba, chiều cao giảm đi một nửa thì thể tích hình hộp chữ nhật đó thay đổi thế nào?
\choice
{Tăng lên $2$ lần}
{\True Tăng lên $3$ lần}
{Tăng lên $6$ lần}
{Không đổi}
\loigiai{
Thể tích ban đầu: $V = a \times b \times c$.\\
Thể tích mới:
$$V' = (2 \times a) \times (3 \times b) \times \left(\dfrac{1}{2} \times c\right) = (2 \times 3 \times \dfrac{1}{2}) \times (a \times b \times c) = 3 \times V.$$
Vậy thể tích tăng lên $3$ lần.\\
\textbf{Đáp án của câu hỏi là B}
}
\end{ex}

\begin{ex} %[5H9V52]
Một bể nước hình hộp chữ nhật có chiều dài $2\text{ m}$, chiều rộng $1\text{ m}$, hiện đang chứa $1\,200\text{ lít}$ nước. Người ta dùng hết $400\text{ lít}$ nước thì mực nước trong bể giảm đi bao nhiêu xăng-ti-mét?
\choice
{$15\text{ cm}$}
{\True $20\text{ cm}$}
{$25\text{ cm}$}
{$30\text{ cm}$}
\loigiai{
Diện tích đáy của bể nước là:
$$2 \times 1 = 2\text{ (m}^2\text{)} = 200\text{ dm}^2.$$
Đổi $400\text{ lít} = 400\text{ dm}^3$.\\
Chiều cao mực nước giảm đi là:
$$400 : 200 = 2\text{ (dm)} = 20\text{ (cm)}.$$
\textbf{Đáp án của câu hỏi là B}
}
\end{ex}

\begin{ex} %[5H9V52]
Một khối đá hình hộp chữ nhật có diện tích xung quanh là $30\text{ dm}^2$, chiều cao $2{,}5\text{ dm}$. Biết chiều dài hơn chiều rộng $2\text{ dm}$. Thể tích của khối đá đó là
\choice
{$18\text{ dm}^3$}
{\True $20\text{ dm}^3$}
{$24\text{ dm}^3$}
{$32\text{ dm}^3$}
\loigiai{
Chu vi đáy của khối đá là:
$$30 : 2{,}5 = 12\text{ (dm)}.$$
Nửa chu vi đáy (tổng chiều dài và chiều rộng) là:
$$12 : 2 = 6\text{ (dm)}.$$
Chiều dài đáy là:
$$(6 + 2) : 2 = 4\text{ (dm)}.$$
Chiều rộng đáy là:
$$6 - 4 = 2\text{ (dm)}.$$
Thể tích khối đá là:
$$4 \times 2 \times 2{,}5 = 20\text{ (dm}^3\text{)}.$$
\textbf{Đáp án của câu hỏi là B}
}
\end{ex}

\begin{ex} %[5H9V52]
Một bể cá hình hộp chữ nhật có chiều dài $60\text{ cm}$, rộng $40\text{ cm}$ và cao $40\text{ cm}$. Mực nước ban đầu cao $25\text{ cm}$. Người ta thả vào bể $5$ viên bi sắt đặc giống nhau thì thấy mực nước dâng lên đến $28\text{ cm}$. Thể tích của mỗi viên bi sắt là
\choice
{$1\,200\text{ cm}^3$}
{\True $1\,440\text{ cm}^3$}
{$1\,500\text{ cm}^3$}
{$7\,200\text{ cm}^3$}
\loigiai{
Mực nước dâng lên thêm là:
$$28 - 25 = 3\text{ (cm)}.$$
Tổng thể tích của $5$ viên bi sắt chính là lượng nước dâng lên:
$$60 \times 40 \times 3 = 7\,200\text{ (cm}^3\text{)}.$$
Thể tích của mỗi viên bi sắt là:
$$7\,200 : 5 = 1\,440\text{ (cm}^3\text{)}.$$
\textbf{Đáp án của câu hỏi là B}
}
\end{ex}

\begin{ex} %[5H9V52]
Người ta đào một cái hố dạng hình hộp chữ nhật có chiều dài $5\text{ m}$, rộng $3\text{ m}$ và sâu $1{,}6\text{ m}$. Đất đào lên tơi xốp nên thể tích tăng thêm $15\%$. Thể tích đất thực tế phải vận chuyển đi là
\choice
{$24\text{ m}^3$}
{\True $27{,}6\text{ m}^3$}
{$26{,}4\text{ m}^3$}
{$28\text{ m}^3$}
\loigiai{
Thể tích đất trong hố ban đầu là:
$$5 \times 3 \times 1{,}6 = 24\text{ (m}^3\text{)}.$$
Thể tích đất tơi xốp cần vận chuyển là:
$$24 \times (100\% + 15\%) = 24 \times 1{,}15 = 27{,}6\text{ (m}^3\text{)}.$$
\textbf{Đáp án của câu hỏi là B}
}
\end{ex}

\textbf{Phần 2. Trắc nghiệm đúng sai}
\setcounterpageref{ex}{0}

\begin{ex} %[5H9V52]
Một bể nước hình hộp chữ nhật có các kích thước trong lòng bể là: dài $2{,}5\text{ m}$, rộng $2\text{ m}$, cao $1{,}6\text{ m}$. Hiện tại bể đang chứa $5\,000\text{ lít}$ nước. Xét tính đúng/sai của các nhận định sau:
\choiceTF
{\True Dung tích toàn bộ của bể là $8\text{ m}^3$}
{\True Lượng nước hiện tại đang chiếm $62{,}5\%$ dung tích bể}
{Mực nước hiện tại trong bể cao $1{,}2\text{ m}$}
{\True Cần bơm thêm $3\,000\text{ dm}^3$ nước nữa để bể đầy}
\loigiai{
\begin{itemchoice}
\item Đúng. Dung tích bể: $2{,}5 \times 2 \times 1{,}6 = 8\text{ (m}^3\text{)} = 8\,000\text{ lít}$.
\item Đúng. Tỉ số phần trăm: $5\,000 : 8\,000 = 62{,}5\%$.
\item Sai. Diện tích đáy bể là $2{,}5 \times 2 = 5\text{ (m}^2\text{)}$. Đổi $5\,000\text{ lít} = 5\text{ m}^3$. Mực nước cao $5 : 5 = 1\text{ (m)}$ (chứ không phải $1{,}2\text{ m}$).
\item Đúng. Cần bơm thêm: $8\,000 - 5\,000 = 3\,000\text{ lít} = 3\,000\text{ dm}^3$.
\end{itemchoice}
}
\end{ex}

\begin{ex} %[5H9V52]
Có hai hình hộp chữ nhật: Hình $A$ có kích thước $8\text{ cm} \times 6\text{ cm} \times 5\text{ cm}$ và Hình $B$ có kích thước $10\text{ cm} \times 4\text{ cm} \times 6\text{ cm}$. Xét tính đúng/sai của các nhận định sau:
\choiceTF
{\True Thể tích của hình $A$ và thể tích của hình $B$ bằng nhau}
{\True Thể tích của mỗi hình bằng $240\text{ cm}^3$}
{Diện tích toàn phần của hình $A$ nhỏ hơn diện tích toàn phần của hình $B$}
{\True Diện tích đáy của hình $A$ lớn hơn diện tích đáy của hình $B$}
\loigiai{
\begin{itemchoice}
\item Đúng. $V_A = 8 \times 6 \times 5 = 240\text{ cm}^3$, $V_B = 10 \times 4 \times 6 = 240\text{ cm}^3$.
\item Đúng. Cả hai đều bằng $240\text{ cm}^3$.
\item Sai. $S_{tp}(A) = 2 \times (8 \times 6 + 8 \times 5 + 6 \times 5) = 2 \times (48 + 40 + 30) = 236\text{ cm}^2$.\\
$S_{tp}(B) = 2 \times (10 \times 4 + 10 \times 6 + 4 \times 6) = 2 \times (40 + 60 + 24) = 248\text{ cm}^2$.\\
Vậy $S_{tp}(A) < S_{tp}(B)$ là đúng, nhận định "nhỏ hơn" là đúng. Sửa lại: Ta phát biểu nhận định sai: "Diện tích toàn phần của hình $A$ lớn hơn diện tích toàn phần của hình $B$". Sửa câu c thành "Diện tích toàn phần của hình A lớn hơn diện tích toàn phần của hình B" là Sai.
\item Đúng. Đáy hình $A$ có diện tích $8 \times 6 = 48\text{ cm}^2 > 10 \times 4 = 40\text{ cm}^2$ của hình $B$.
\end{itemchoice}
}
\end{ex}

\textbf{Phần 3. Câu trả lời ngắn}
\setcounterpageref{ex}{0}

\begin{ex} %[5H9V52]
Một bể nước hình hộp chữ nhật có diện tích đáy là $6\text{ m}^2$. Người ta mở vòi cho nước chảy vào bể được $7\,200\text{ lít}$ nước. Hỏi mực nước trong bể cao bao nhiêu đề-xi-mét?
\shortans{12}
\loigiai{
Đổi: $6\text{ m}^2 = 600\text{ dm}^2$.\\
Đổi: $7\,200\text{ lít} = 7\,200\text{ dm}^3$.\\
Chiều cao mực nước là:
$$7\,200 : 600 = 12\text{ (dm)}.$$
}
\end{ex}

\begin{ex} %[5H9V52]
Một khối gỗ hình hộp chữ nhật dài $12\text{ dm}$, rộng $8\text{ dm}$ và cao $5\text{ dm}$. Người ta muốn cắt khối gỗ thành các khối lập phương cạnh $2\text{ dm}$. Hỏi cắt được nhiều nhất bao nhiêu khối lập phương nhỏ như vậy (không tính hao hụt khi cưa)?
\shortans{60}
\loigiai{
Số khối theo chiều dài: $12 : 2 = 6$ (khối).\\
Số khối theo chiều rộng: $8 : 2 = 4$ (khối).\\
Số khối theo chiều cao: $5 : 2 = 2$ (dư $1\text{ dm}$).\\
Số khối lập phương nhỏ cắt được nhiều nhất là:
$$6 \times 4 \times 2 = 60\text{ (khối)}.$$
}
\end{ex}

\begin{ex} %[5H9V52]
Một bể cá hình hộp chữ nhật dài $90\text{ cm}$, rộng $50\text{ cm}$ đang chứa nước cao $35\text{ cm}$. Người ta thả vào bể một hòn đá có thể tích $45\text{ dm}^3$. Hỏi mực nước dâng lên thêm bao nhiêu xăng-ti-mét (nước không tràn ra ngoài)?
\shortans{10}
\loigiai{
Đổi: $45\text{ dm}^3 = 45\,000\text{ cm}^3$.\\
Diện tích đáy của bể cá là:
$$90 \times 50 = 4\,500\text{ (cm}^2\text{)}.$$
Chiều cao mực nước dâng thêm là:
$$45\,000 : 4\,500 = 10\text{ (cm)}.$$
}
\end{ex}

\begin{ex} %[5H9V52]
Một bể bơi mini hình hộp chữ nhật dài $15\text{ m}$, rộng $6\text{ m}$, sâu $1{,}4\text{ m}$. Hiện tại bể đã có $63\text{ m}^3$ nước. Cần bơm thêm vào bể bao nhiêu mét khối nước nữa để đầy bể?
\shortans{63}
\loigiai{
Thể tích bể bơi là:
$$15 \times 6 \times 1{,}4 = 126\text{ (m}^3\text{)}.$$
Số mét khối nước cần bơm thêm là:
$$126 - 63 = 63\text{ (m}^3\text{)}.$$
}
\end{ex}

\begin{ex} %[5H9V52]
Một hình hộp chữ nhật có diện tích xung quanh là $120\text{ cm}^2$, chiều cao là $5\text{ cm}$. Biết chiều dài gấp đôi chiều rộng. Thể tích của hình hộp chữ nhật đó là bao nhiêu xăng-ti-mét khối?
\shortans{160}
\loigiai{
Chu vi đáy của hình hộp là:
$$120 : 5 = 24\text{ (cm)}.$$
Nửa chu vi đáy là:
$$24 : 2 = 12\text{ (cm)}.$$
Chiều rộng đáy là:
$$12 : (2 + 1) = 4\text{ (cm)}.$$
Chiều dài đáy là:
$$4 \times 2 = 8\text{ (cm)}.$$
Thể tích của hình hộp chữ nhật là:
$$8 \times 4 \times 5 = 160\text{ (cm}^3\text{)}.$$
}
\end{ex}

\begin{ex} %[5H9V52]
Một chiếc xe tải chở một khối cát hình hộp chữ nhật có kích thước dài $3\text{ m}$, rộng $2\text{ m}$, cao $0{,}75\text{ m}$. Biết mỗi mét khối cát cân nặng $1\,400\text{ kg}$. Hỏi toàn bộ khối cát đó cân nặng bao nhiêu ki-lô-gam?
\shortans{6300}
\loigiai{
Thể tích của khối cát là:
$$3 \times 2 \times 0{,}75 = 4{,}5\text{ (m}^3\text{)}.$$
Khối lượng cát là:
$$4{,}5 \times 1\,400 = 6\,300\text{ (kg)}.$$
}
\end{ex}
